\section{Лінійна регресія}

У математичній статистиці лінійна регресія — це метод побудови залежності між змінною
y та векторною залежною змінною X.

При використанні цього алгоритму машинного навчання взаємозв'язок між даними будується
за допомогою лінійних функцій, причому невідомі параметри функції отримуються за
вхідними даними. Так само, як і інші регресійні моделі, лінійна регресія повертає
розподіл умовної ймовірності y залежно від X.

Для знаходження параметрів лінійної регресії загалом використовується метод
найменших квадратів (МНК). Також існує можливість використання методу найменших квадратів
для нелінійних моделей.

Якщо говорити мовою формул \eqref{eq:formula-5.1}:
\begin{equation}
	y = f(x, b) + \varepsilon, \quad E(x) = 0 \label{eq:formula-5.1}
\end{equation}
де $b$ --- параметри моделі, $\varepsilon$ --- випадкова похибка моделі.

називається лінійною регресією, якщо $f(x, b)$ подається у вигляді \eqref{eq:formula-5.2}:
\begin{equation}
	f(x, b) = b_{1} x_{1} + b_{2} x_{2} + \dots + b_{k} x_{k} = \sum_{j=1}^{k}b_{j}
	x_{j} = x^{T} \label{eq:formula-5.2}
\end{equation}
$b_{j}$ --- параметри регресії, $x_{j}$ --- регресори, $k$ --- кількість факторів
у моделі

Також виділяють і так звану множинну регресію, що подається у вигляді \eqref{eq:formula-5.3}:
\begin{equation}
	Y = b_{0} + b_{1} x_{i1}+ \dots + b_{j} x_{ij}+ \dots + b_{k} x_{ik}+ e_{i} \label{eq:formula-5.3}
\end{equation}

Вищеописане звісно є частковим випадком того, що можна визначити за кількості вимірювань,
щоб однозначно визначити якого вигляду набуває параметр β.

Використовуючи саме алгоритм баєсівської регресії, можна отримати ще додаткову
інформацію, а саме апріорну ймовірність розподілу. Ці ймовірності сильно взаємопов'язані
із функцією правдоподібності даних відповідно до теореми Баєса для отримання вже
апостеріорних ймовірностей важливих параметрів βтаσ².

Ці ймовірності дуже залежать від кількості та якості інформації, що подається на
вхід базової моделі, оскільки функція лінійна.

Щодо версії алгоритму машинного навчання, яка застосовувалась у роботі, то це Bayesian
ridge, що використовує регуляризацію з нормою L2, тоді як класична баєсова регресія
--- це регресійна модель, визначена з явними апріорними ймовірностями та
параметрами. Вибір ймовірностей може мати ефект регуляризації: наприклад, використання
розподілу Лапласа для коефіцієнтів еквівалентно регуляризації L1. Вони не
однакові, тому що Bayesian ridge є своєрідною моделлю регресії, а баєсівський
підхід є загальним способом визначення та оцінки статистичних моделей.
